\begin{frame}{$\phi$ 는 위치만, 기울기는 그대로}
  \begin{columns}[T]
    \begin{column}{0.4\textwidth}
      \vskip0.4em
      \begin{itemize}
        \item \kb{사전확률 $\phi$ 는 방향이 아니라 위치($b$)만
          움직인다.}
        \item $\phi$ 는 $\log\frac{\phi}{1-\phi}$ 항에만 등장 ---
          클래스 비율이 무엇이든 \textbf{직선의 기울기는 그대로},
          \textbf{밀리는 정도}만 바뀜.
      \end{itemize}
    \end{column}
    \begin{column}{0.58\textwidth}
      \centering
      \includegraphics[height=5.0cm]{ch03_prior_shift.png}
    \end{column}
  \end{columns}
\note{\begin{itemize}
\item 다음 ``가정이 깨질 때'' 절에서 이 효과를 \textbf{숫자로}
      확인.
\item {\footnotesize (전개 과정의 전체 유도 = 이번 장 연습문제 2번.
  ``정확성 확인''에서 $\Sigma_0\ne\Sigma_1$ 이면 이차항 미상쇄
  $\to$ 곡선 = QDA.)}
\end{itemize}}
\end{frame}
